\scanpage{20}
% L1–2
$C$ un $\mathcal U$-site, $E$ catégorie où les petites $\mathop{\underline{\lim}}$ sont repr. déterminer les foncteurs

% L3
\[
u\colon\widetilde C\longrightarrow E
\]

% L4–5
commutant aux $\mathop{\underline{\lim}}$ (petites) ? Connu \inserted{à isom près} quand on connait $u\circ\varepsilon_C=u_0$

% L6
\[
u_0\colon C\longrightarrow E
\]

% L6-right
\diagram{\begin{tikzcd}[ampersand replacement=\&]
C\arrow[r,"{\varepsilon_C}"]\arrow[rr,bend right=27,"{u_0}"'] \& \widetilde C\arrow[r,"u"] \& E
\end{tikzcd}}

% L7
car pour $F\in\widetilde C$, on a un isom. fonctoriel en $F$

% L8
\[
F=\varinjlim_{C/F}\varepsilon_C(x)
\]

% L9
\[
\text{d'où}\qquad u(F)=\mathop{\underline{\lim}}_{C/F}u_0(X).
\]

% L10
De façon précise, on trouve un foncteur $u\longmapsto u\circ\varepsilon_C$

% L11
\[
\alpha\colon\underline{\mathrm{Hom}}'(\widetilde C,E)\longrightarrow\underline{\mathrm{Hom}}\,(C,E)
\]

% L12–14
qui est \uline{pl. fid}, le signe $\underline{\mathrm{Hom}}'$ désignant les foncteurs commutant aux petites $\varinjlim$. En fait, on a un foncteur ``prolongement à $\widetilde C$''

% L14
\[
\beta\colon\underline{\mathrm{Hom}}\,(C,E)\longrightarrow\underline{\mathrm{Hom}}\,(C,E)
\]

% L15
défini par

% L16
\[
\beta(u_0)=F\longmapsto\mathop{\underline{\lim}}_{C/F}u_0(X)
\]

% L17
et on a un isom canonique

% L18
\[
\beta\unsure{\circ|\xi}\alpha\simeq i\quad
(\text{inclusion }\deleted{\unsure{\text{de}}}\ \underline{\mathrm{Hom}}'(\widetilde C,E)\hookrightarrow\underline{\mathrm{Hom}}\,(\unsure{\overline C|\widetilde C},E))
\]

% L19–21
Mais il n'est pas clair que $\forall u_0\colon C\longrightarrow E$, $\beta(u_0)$ soit dans $\underline{\mathrm{Hom}}'$. A vrai dire, $\beta(u_0)$ est induit par $\widehat\beta(u_0)\colon\widehat C\longrightarrow E$ défini de la $=$ façon, et on sait que $\widehat\beta(u_0)$ commute aux $\varinjlim$, en fait,

% L22
\diagram{\begin{tikzcd}[ampersand replacement=\&,column sep=large]
\underline{\mathrm{Hom}}\,(C,E)\arrow[r,shift left,"{\cong}","{\widehat\beta}" description,near start] \& \underline{\mathrm{Hom}}'(\widehat C,E)\arrow[l,shift left,"{\widehat\alpha}"]
\end{tikzcd}}

% L22–24
(cf rappel plus bas). Identifions dorénavant (modulo équivalence) $\underline{\mathrm{Hom}}\,(C,E)$ à $\underline{\mathrm{Hom}}'(\widehat C,E)$, alors $\alpha$ et $\beta$ deviennent des foncteurs $\alpha'$, $\beta'$

% L25–26
\[
\begin{aligned}
\underline{\mathrm{Hom}}'(\widetilde C,E)&\xrightarrow{\alpha'}\underline{\mathrm{Hom}}'(\widehat C,E)\\
\underline{\mathrm{Hom}}'(\widehat C,E)&\xrightarrow{\beta'}\underline{\mathrm{Hom}}\,(\widetilde C,E)
\end{aligned}
\]

% foot
% archive sheet number: 19
